\documentclass[11pt,fontset=fandol]{ctexart}
\usepackage[a4paper,margin=23mm]{geometry}
\usepackage{amsmath,amssymb,mathtools,bm,mathrsfs,esint}
\usepackage{graphicx,longtable,booktabs,array,calc}
\usepackage[normalem]{ulem}
\usepackage{xcolor,hyperref}
\hypersetup{hidelinks}
\setcounter{secnumdepth}{-1}
\newcounter{none}
\providecommand{\tightlist}{\setlength{\itemsep}{0pt}\setlength{\parskip}{0pt}}
\providecommand{\passthrough}[1]{#1}
\usepackage{amssymb,mathrsfs,esint,bm,newunicodechar}
\allowdisplaybreaks
\setlength{\emergencystretch}{3em}
\providecommand{\llangle}{\langle\!\langle}
\providecommand{\rrangle}{\rangle\!\rangle}
\newunicodechar{∞}{\ensuremath{\infty}}
\newunicodechar{∈}{\ensuremath{\in}}
\newunicodechar{∑}{\ensuremath{\sum}}
\newunicodechar{∫}{\ensuremath{\int}}
\newunicodechar{≤}{\ensuremath{\leq}}
\newunicodechar{≥}{\ensuremath{\geq}}
\newunicodechar{≠}{\ensuremath{\neq}}
\newunicodechar{→}{\ensuremath{\to}}
\newunicodechar{⇒}{\ensuremath{\Rightarrow}}
\newunicodechar{α}{\ensuremath{\alpha}}
\newunicodechar{β}{\ensuremath{\beta}}
\newunicodechar{γ}{\ensuremath{\gamma}}
\newunicodechar{λ}{\ensuremath{\lambda}}
\newunicodechar{φ}{\ensuremath{\varphi}}
\newunicodechar{ε}{\ensuremath{\varepsilon}}
\newunicodechar{σ}{\ensuremath{\sigma}}
\newunicodechar{θ}{\ensuremath{\theta}}
\newunicodechar{δ}{\ensuremath{\delta}}
\newunicodechar{η}{\ensuremath{\eta}}
\newunicodechar{π}{\ensuremath{\pi}}
\newunicodechar{ξ}{\ensuremath{\xi}}
\newunicodechar{Σ}{\ensuremath{\Sigma}}
\newunicodechar{Δ}{\ensuremath{\Delta}}
\newunicodechar{∂}{\ensuremath{\partial}}
\newunicodechar{∀}{\ensuremath{\forall}}
\newunicodechar{∃}{\ensuremath{\exists}}
\newunicodechar{√}{\ensuremath{\surd}}
\newunicodechar{−}{\ensuremath{-}}
\newunicodechar{μ}{\ensuremath{\mu}}
\newunicodechar{ρ}{\ensuremath{\rho}}
\newunicodechar{τ}{\ensuremath{\tau}}
\newunicodechar{ω}{\ensuremath{\omega}}
\newunicodechar{Ω}{\ensuremath{\Omega}}
\newunicodechar{χ}{\ensuremath{\chi}}
\newunicodechar{ψ}{\ensuremath{\psi}}
\newunicodechar{Γ}{\ensuremath{\Gamma}}
\newunicodechar{Φ}{\ensuremath{\Phi}}
\newunicodechar{Λ}{\ensuremath{\Lambda}}
\newunicodechar{ℝ}{\ensuremath{\mathbb{R}}}
\newunicodechar{ℤ}{\ensuremath{\mathbb{Z}}}
\newunicodechar{ℕ}{\ensuremath{\mathbb{N}}}
\newunicodechar{△}{\ensuremath{\triangle}}
\newunicodechar{⊗}{\ensuremath{\otimes}}
\begin{document}
\section{2026年杭州电子科技大学601数学分析考研真题}\label{ux5e74ux676dux5ddeux7535ux5b50ux79d1ux6280ux5927ux5b66601ux6570ux5b66ux5206ux6790ux8003ux7814ux771fux9898}

来源：\texttt{HDU/601数学分析/2026/2026年杭州电子科技大学601数学分析考研真题.pdf}。原
PDF 页序：1, 2。

\begin{enumerate}
\def\labelenumi{\arabic{enumi}.}
\setcounter{enumi}{2}
\tightlist
\item
  杭州电子科技大学2026年数学分析
\end{enumerate}

一、填空题（每题5分，共40分）

\begin{enumerate}
\def\labelenumi{\arabic{enumi}.}
\item
  \(\lim_{n \to \infty} \tan^n \left( \frac{\pi}{4} + \frac{1}{n} \right) = \_\_\_\_\_\)
\item
  \(\lim_{(x,y) \to (+\infty,+\infty)} \left( x^2 + y^2 \right) e^{-(x+y)} = \_\_\_\_\_\)
\item
  \(\int_0^{\ln 2} \sqrt{e^x - 1} dx = \_\_\_\_\_\)
\item
  \(\lim_{n \to \infty} \left[ \frac{f\left( 2026 + \frac{1}{n} \right)}{f(2026)} \right]^n = \_\_\_\_\_\)
\item
  已知 \(f(x,y,z) = xy^2 z^3 e^{xyz}\), \(z = z(x,y)\) 由
  \(x+y+z+xyz=0\) 确定, 求 \(f_x(1,1,-1) = \_\_\_\_\_\)
\item
  设 \(f(0) = 0\),
  \(F(t) = \iiint_V \left[ z^2 + f(x^2 + y^2) \right] dxdydz\)
  (\(V: x^2 + y^2 \leq t^2, 0 \leq z \leq 1\)), 求
  \(\lim_{t \to 0^+} \frac{F(t)}{t^2} = \_\_\_\_\_\)
\item
  \(\int_L 3x^2 y dx + (x^3 + x - 2y) dy = \_\_\_\_\_\), 其中 \(L\) 为沿
  \(x^2 + y^2 = 2x\) 从 \((0,0) \to (2,0)\), 再沿 \(x^2 + y^2 = 4\) 从
  \((2,0) \to (0,2)\) 的曲线.
\item
  \(\int_0^{+\infty} \frac{\ln x}{x^2 + a} dx = \_\_\_\_\_\)
\end{enumerate}

二、计算题（每题7分，共28分）

\begin{enumerate}
\def\labelenumi{\arabic{enumi}.}
\item
  \(\lim_{x \to +\infty} \left( \frac{a^x - 1}{(a-1)x} \right)^{\frac{1}{x}}\)
  (\(a > 0, a \neq 1\)).
\item
  \(\int_{\frac{1}{2}}^{\frac{3}{4}} \frac{\arcsin \sqrt{x}}{\sqrt{x(1-x)}} dx\).
\item
  已知
  \(f(x) = \begin{cases} \frac{\sin x}{x}, & x \neq 0 \\ 1, & x = 0 \end{cases}\),
  求 \(f^{(k)}(0)\) (\(k\) 阶导数).
\item
  \(I = \iint_S f(x,y,z) dS\), 其中 \(S: x^2 + y^2 + z^2 = 1\),
  \(f(x,y,z) = \begin{cases} \sqrt{x^2 + y^2}, & 0 \leq z \leq \sqrt{x^2 + y^2} \\ 0, & z < 0 \text{ 或 } z > \sqrt{x^2 + y^2} \end{cases}\)
\end{enumerate}

三、综合题（每题10分，共40分）

\begin{enumerate}
\def\labelenumi{\arabic{enumi}.}
\tightlist
\item
  已知 \(y = y(x)\) 由 \(x = \int_1^{y-x} e^{-u^2} du\) 确定, \(y\)
  二阶可导, 求 \(\lim_{x \to 0} \frac{y - (1+e)x - 1}{x^2}\).
\end{enumerate}

2、设
\(a_{n}=\frac{1}{2\ln 2}+\frac{1}{3\ln 3}+\cdots+\frac{1}{n\ln n}-\ln \ln n\)
，判断 \(a_{n}\) 是否收敛. 3、设 \(u=f(z)\) ， \(z=x+y\varphi(z)\)
，证明：
\(\frac{\partial^{n}u}{\partial y^{n}}=\frac{\partial^{n-1}}{\partial x^{n-1}}\left[\varphi(z)^{n}\frac{\partial u}{\partial x}\right].\)
4、设 \(I(\alpha)=\int_{0}^{+\infty}\frac{e^{-x}}{|\sin x|^{\alpha}}dx\)
，证明 \(I(\alpha)\) 对 \(\alpha \in[0,b](0<b<1)\) 一致收敛.
四、证明题（1-3题10分，4题12分，共42分） 1、已知 \(f(x)\) 一致连续，
\(\varphi(x)\) 连续，且 \(\lim_{x \to +\infty}[f(x)-\varphi(x)]=0\)
，证明 \(\varphi(x)\) 在 \([a,+\infty)\) 上一致连续. 2、已知 \(f(x)\) 在
\([a,b]\) 连续， \((a,b)\) 可微，且 \(f(a)\neq f(b)\) ，证明：
\(\exists \xi,\eta \in(a,b)\) ，使得
\(f^{\prime}(\xi)=\frac{a+b}{2\eta}f^{\prime}(\eta).\) 3、已知 \(f(x)\)
在 \([0,1]\) 连续， \(f_{1}(x)=f(x)\) ，
\(f_{n+1}(x)=\int_{x}^{1}f_{n}(t)dt\) （ \(n\geq 1\) ），证明
\(\sum_{n=1}^{\infty}f_{n}(x)\) 在 \(0\leq x\leq 1\) 上一致收敛. 4、设
\(\varphi(x)\) 在 \(\mathbb{R}\) 上周期为1，且
\(\int_{0}^{1}\varphi(x)dx=0\) ， \(f(x)\) 在 \([0,1]\) 连续可导，
\(a_{n}=\int_{0}^{1}f(x)\varphi(nx)dx\) ，证明级数
\(\sum_{n=1}^{\infty}a_{n}^{2}\) 收敛.
\end{document}
